% Draft Statutory Text. Enactable legislative language in the (a)(1)(i)
% hierarchy, in the style of the template's new Subpart J. The ten-gate
% thresholds are codified in one table sourced from the VVUQ-02 gate spec and
% standards map; the prose references the table rather than restating numbers.

The operative text of H.R. 9510 follows. It is drafted as enactable legislative
language and is preceded by the customary bill header and enacting clause.



\begin{flushleft}
\textbf{119th Congress, 2d Session \quad H. R. 9510}\\[3pt]
\textbf{A BILL}\\[2pt]
To require an automated verification, validation, and uncertainty quantification
process to clear robot-patient interaction code before that code is generated or
executed in a Physical AI oncology clinical trial, and for other purposes.
\end{flushleft}

\clearpage

\textit{Be it enacted by the Senate and House of Representatives of the United
States of America in Congress assembled,}

\subsection{Section 1. Short Title}

(a) \textit{Short title.} This Act may be cited as the ``Verification Before
Generation in Physical AI Oncology Trials Act of 2026''.

(b) \textit{Popular name.} This Act may also be referred to as the VVUQ Physical
AI Oncology Trial Bill.

\subsection{Section 2. Definitions}

(a) \textit{Incorporation by reference.} The terms used in this Act have the
meanings given them in the definitions set out in
Section~\ref{sec:definitions} of this instrument, which are incorporated by
reference and not restated here.

(b) \textit{Operative terms.} The terms on which the operative sections depend
include: Physical AI system; robot-patient interaction; VVUQ gate; verification
fraction; Unification Standard Level (USL); Physical AI Standard Level (PSL);
high-risk artificial intelligence; and AI-enabled device software function.

\subsection{Section 3. Verification Before Generation}

(a) \textit{Requirement.} No robot-patient interaction code may be generated or
executed in a Physical AI oncology clinical trial unless an automated VVUQ
verification process, bound to one or more named external standards, has first
cleared for that interaction, and the cleared verification record has been
documented and attested under Section 6.

(b) \textit{Priority.} The verification process is held to a higher standard
than, and shall precede, the generated code. Where resources or time are
constrained, the verification process takes precedence over the generation step,
because the verification is the decision-bearing work and the generation is
not \cite{kawchak2026vvuq01, kawchak2026vvuq02}.

(c) \textit{Order of operations.} A sponsor shall not generate, and shall not
execute, robot-patient interaction code in advance of the cleared verification
record. A record created after generation or execution does not satisfy this
section.

(d) \textit{Scope.} For purposes of this section, robot-patient interaction code
is any algorithm or generated software that governs a Physical AI system's
physical contact with, or action upon, a human subject, whether the system is
classified as a medical product, a research tool, or both. The verification
process and the credibility framework it serves track the agency credibility
guidance and the ASME V\&V 40 standard \cite{asmevv40, fda2023credibility}.

\subsection{Section 4. Validation Requirements and Gate Thresholds}

(a) \textit{Verification fraction.} For each robot-patient interaction, the
verification fraction shall equal 1.0; every verification check shall pass. This
is an equality, not an inequality with slack.

(b) \textit{Validation agreement.} The result shall agree with an independent
reference at or above the per-gate agreement threshold, with the maximum relative
error at or below the per-gate bound, and with the human review recorded.

(c) \textit{Uncertainty bound.} The coefficient of variation across at least
three seeded runs shall not exceed the per-gate bound; divergence beyond the
bound sets ESCALATE.

(d) \textit{Hard predicate.} For immediate-catastrophe interactions, an
additional pass/fail predicate shall also hold: zero breaches of a vascular
no-fly volume, a minimum human clearance in a shared operating room, and a
guaranteed safe state under fault or emergency stop. These gates carry a
verification fraction of 1.0 and the tightest dispersion bounds
\cite{iso14971, isots15066}.

(e) \textit{Decision rule.} The gate emits ACCEPT, BLOCK, or ESCALATE. Any single
failing dimension blocks; ambiguity or divergence escalates and hands the
interaction back to a qualified human.

(f) \textit{Threshold schedule.} The per-gate agreement and dispersion bounds,
the gates that carry a hard predicate, and the external standard that governs
each gate are codified in Table~\ref{tab:statute-thresholds}; the numeric values
are not restated in this subsection \cite{nasastd7009a, iec8060127, ul4600,
ieee7009}.

\begin{xltabular}{\textwidth}{L{5cm} L{1.25cm} L{1.3cm} L{1cm} Y}
\toprule
\textbf{Gate} & \textbf{Agree $\geq$} & \textbf{Max CV $\leq$} & \textbf{Hard pred.} & \textbf{Governing standards} \\
\midrule
\endfirsthead
\toprule
\textbf{Gate} & \textbf{Agree $\geq$} & \textbf{Max CV $\leq$} & \textbf{Hard pred.} & \textbf{Governing standards} \\
\midrule
\endhead
\midrule
\multicolumn{5}{r}{\textit{Continued on next page}} \\
\endfoot
\bottomrule
\endlastfoot
01 Bimanual hand-eye servo & 0.97 & 0.08 & No & IEC 80601-2-77; ISO 9283; ASME V\&V 40 \\
02 Dexterous finger force & 0.95 & 0.10 & No & IEC 80601-2-77; ISO/TS 15066 \\
03 Whole-body balance & 0.98 & 0.06 & No & ISO 13482; IEC 60601-1 \\
04 Autonomous plan correctness & 0.95 & 0.10 & No & UL 4600; IEEE 7009; IEC 62304 \\
05 Instrument grasp and handover & 0.96 & 0.10 & No & IEC 80601-2-77; ISO 9283 \\
06 Vascular no-fly (hand) & 1.00 & 0.05 & Yes & IEC 80601-2-77; ISO 14971 \\
07 Bimanual suturing anastomosis & 0.96 & 0.08 & No & IEC 80601-2-77; Fistula Risk Score \\
08 Perception scene understanding & 0.95 & 0.10 & No & ASME V\&V 40; IEC 62304 \\
09 Shared-OR human collision & 1.00 & 0.06 & Yes & ISO/TS 15066; ISO 10218-1; ISO 13482 \\
10 Fault, e-stop, degradation & 1.00 & 0.05 & Yes & \small{IEC 60601-1/ISO 13849-1/IEEE 7009} \\
\end{xltabular}

\subsection{Section 5. Readiness Gates}

(a) \textit{Site gate.} Before enrollment, the trial site shall pass the Physical
AI Standard Level gate on all three dimensions: legislative authorization,
regulatory compliance, and operational readiness \cite{kawchak2026national}.

(b) \textit{Robot gate.} Before enrollment, each robot shall meet the minimum
Unification Standard Level score for its procedure type: surgical, at least 7.0;
therapeutic positioning and diagnostic needle placement, at least 6.0;
rehabilitative, at least 4.0; and companion monitoring, at least 3.0.

(c) \textit{Phase 0 simulation validation.} Before enrollment, each robot's
behavior shall be validated across at least two simulation frameworks within the
cross-framework tolerances of under 2 mm trajectory deviation and under 0.5 N
force discrepancy. The site gate, the robot gate, and the code-level gate of
Sections 3 and 4 shall all clear before a patient is enrolled or a robot acts.

\subsection{Section 6. Documentation and Attestation}

(a) \textit{Algorithm documentation.} For each covered algorithm, the sponsor
shall file the algorithm documentation specified in
Section~\ref{sec:algorithm_documentation}, including the
verification-before-generation record.

(b) \textit{Compliance statement and attestation.} The sponsor shall file the
artificial-intelligence compliance statement and the signed attestation specified
in Section~\ref{sec:attestations_compliance}, consistent with the state
precedents \cite{tx-sb1822, ny-a9149, ca-sb1120}.

(c) \textit{Filed before generation.} The verification record shall be filed
before the covered code is generated or executed.

\subsection{Section 7. Cybersecurity, Human Oversight, and Lifecycle}

(a) \textit{Cybersecurity.} Each Physical AI system shall incorporate
authentication and access control, encryption of data in transit and at rest,
network segmentation, software-integrity verification, and a hash-chained,
cryptographically signed audit trail, satisfying electronic-records rule
\cite{cfr-part11, kawchak2026mcp}.

(b) \textit{Human oversight.} Each robot-patient interaction shall be conducted
under recorded human oversight, with a manual override and an emergency stop that
brings the system to a safe state within the procedure-appropriate budget, and
with a hand-back-to-human escalation whenever the gate escalates or a fault is
detected \cite{iec60601}.

(c) \textit{Lifecycle and decommissioning.} The sponsor shall maintain
configuration and model-version management, monitor for model and data drift, and,
on decommissioning, archive the audit trail, securely erase patient data, and
revoke access.

\subsection{Section 8. Nondiscrimination}

(a) \textit{Bias minimization.} Each covered algorithm and its training data
shall minimize the risk of bias on protected characteristics and shall adhere to
evidence-based clinical guidelines \cite{usc-titlevi, ny-a9149}.

(b) \textit{Accessibility.} Each Physical AI system shall be designed and operated
consistent with the accessibility and accommodation requirements of the Americans
with Disabilities Act \cite{usc-ada}.

\subsection{Section 9. Enforcement and Effective Date}

(a) \textit{Enforcement.} The oversight roles, audit and penalty mechanics, and
the operative trigger, which is the absence of a cleared
verification-before-generation record, are set out in
Section~\ref{sec:implementation_enforcement} and are not restated here.

(b) \textit{Effective date.} The effective date and the phasing of the readiness
and documentation duties are set out in Section~\ref{sec:conclusions} and are not
restated here.

\vspace{0.4em}
\noindent\textit{Authority and source.} This Act adapts material in the public
domain under 17 U.S.C. § 105, including 21 CFR Parts 312 and 50, and is not
endorsed or sponsored by CFR, ICH, or FDA. The validation thresholds are drawn
from the author's second-study gate specification and standards map, and the
new-subpart structure follows the prior 21 CFR Part 312 adaptation.
